\chapter{Modelos de inventarios}
\label{ch:inventarios}

\section{Qué decisión se modela}
Una política de inventario combina cuánto reponer y cuándo hacerlo. El costo
mínimo depende de las unidades de tiempo, del patrón de demanda y del costo de
faltantes; no basta con introducir una demanda media
\cite{harris1913many,silver1998inventory}. QSB distingue EOQ/EPQ, descuentos
por cantidad, vendedor de periódicos, lotes por períodos y cuatro políticas de
revisión estocástica. Todos utilizan el sobre JSON \texttt{kind}/\texttt{model}
y el backend nativo o de validación. No hay backend externo de inventarios.

\section{EOQ y producción gradual}
Con demanda constante \(D\), costo fijo \(S\) por pedido y mantenimiento \(H\)
por unidad y período, sin faltantes, el costo relevante es
\[
 C(Q)=DS/Q+HQ/2,\qquad Q^*=\sqrt{2DS/H}.
\]
La adquisición \(Dc\) se suma al costo total, pero no altera el óptimo si el
precio unitario \(c\) es constante. El plazo \(\ell\) produce el punto de
reorden \(r=D\ell\). Tasas, plazo y mantenimiento deben usar el mismo período.
Para producción gradual a tasa \(p>D\), el inventario medio es
\(Q(1-D/p)/2\) y
\[
 Q^*_{EPQ}=\sqrt{\frac{2DS}{H(1-D/p)}}.
\]
El cálculo es continuo: comprar cajas o unidades enteras requiere comparar
cantidades factibles. QSB no introduce automáticamente restricciones de embalaje.

\subsection{Ejemplo: reposición de cuadernos}
El archivo \path{manual/examples/prob_eoq.json} contiene \(D=1200\)
unidades/año, \(S=30\), \(H=2\), \(c=5\), plazo \(0.01\) años y cantidad
actual de pedido \(100\). Desde la raíz del repositorio:
\begin{lstlisting}
.build/debug/qsb validate-inventory-json manual/examples/prob_eoq.json
.build/debug/qsb solve-inventory-json manual/examples/prob_eoq.json
\end{lstlisting}
Se obtiene \(Q^*=189.736660\), \(6.324555\) pedidos/año, ciclo
\(0.158114\) años y punto de reorden \(12\). Ordenar y mantener cuestan
\(189.736660\) cada uno; el costo relevante es \(379.473319\) y el total
\(6379.473319\). Pedir \(100\) cuesta \(6460\): el ahorro anual del modelo
es \(80.526681\), no el valor completo de las compras. El documento de salida
incluye la comparación en \texttt{knownQuantity} y exactitud
\texttt{closedForm}. En \path{prob_epq.json}, \(p=2400\) eleva el lote
óptimo a \(268.328157\), con costo relevante \(268.328157\).

\section{Descuentos sobre todas las unidades}
QSB evalúa cada tramo de precio y sus cantidades factibles, incluyendo
adquisición, preparación y mantenimiento. El descuento se aplica a todas las
unidades del pedido; no es una tarifa marginal por bloques. El mantenimiento
es un costo por unidad, no un porcentaje recalculado a partir del precio.
\path{prob_descuento.json} añade un descuento del \(10\%\) a partir de
\(300\) unidades. Con \texttt{solve-inventory-json}, resulta elegido
\(Q=300\), precio \(4.5\), costo de preparación \(120\), mantenimiento
\(300\), adquisición \(5400\) y total \(5820\). El costo relevante aumenta a
\(420\), pero el ahorro en adquisición hace preferible el tramo. La tabla de
candidatos y \texttt{optimum} explican la decisión; la enumeración de tramos
es exacta dentro del modelo continuo.

\section{Vendedor de periódicos}
Se decide una sola reposición para un horizonte sin recuperación posterior.
Para precio \(v\), adquisición \(c\), salvamento \(s\) y penalización
adicional por faltante \(b\), los costos marginales son
\(C_u=v-c+b\) y \(C_o=c-s\). La fracción crítica es
\[
 F(q^*)=\frac{C_u}{C_u+C_o}.
\]
QSB admite demanda normal, calcula su cuantil y los valores esperados de ventas,
sobrantes, faltantes y beneficio. La probabilidad de cubrir toda la demanda es
el nivel de servicio del ciclo; no es la fracción esperada de demanda servida.
En \path{prob_newsboy.json}, media \(100\), desviación \(20\),
\(v=5,c=2,s=1,b=1\) y costo fijo cero dan fracción crítica \(0.8\), pedido
\(116.832425\), ventas esperadas \(97.767247\), sobrantes \(19.065178\),
faltantes \(2.232753\) y beneficio \(272.003808\). El cuantil normal utiliza
aritmética aproximada; la etiqueta \texttt{closedForm} describe el método,
no una cantidad decimal exacta ni una garantía para demanda no normal.

\section{Dimensionamiento de lotes}
El inventario final satisface \(I_t=I_{t-1}+q_t-d_t\). QSB minimiza
\[
 \sum_t\left(S_t\mathbf{1}_{q_t>0}+c_tq_t+
 h_t\max(I_t,0)+b_t\max(-I_t,0)\right).
\]
La programación dinámica usa demanda y cantidades enteras, inventario inicial
cero y final cero, permitiendo atraso con costo. No representa capacidad de
producción ni stock inicial configurable. La exactitud se refiere al espacio
de estados acotado de esta implementación educativa: inventarios entre
menos y más la demanda total, y producción por período entre cero y esa demanda.
El costo computacional crece con la demanda total y el horizonte.
\path{prob_lotes.json} tiene tres demandas de diez unidades, preparación
\(20\), producción unitaria \(2\), mantenimiento \(1\) y atraso \(10\).
QSB devuelve producción \((20,0,10)\), inventario final \((10,0,0)\) y costo
\(110\). La política \((10,20,0)\) también cuesta \(110\); un resultado
óptimo no implica unicidad. La tabla de decisiones permite comprobar el balance
y la imputación de costos a cada período.

\section{Revisión con demanda incierta}
Las cuatro políticas usan demanda normal por período y plazo constante.
Con incremento independiente de demanda, el período de protección \(T\)
tiene media \(DT\) y desviación \(\sigma\sqrt T\). El stock de seguridad
se expresa mediante el cuantil normal y el costo esperado de faltantes usa la
función de pérdida normal. Son aproximaciones de cantidades continuas,
marcadas \texttt{approximate}; no resuelven una política óptima general para
cualquier distribución o plazo aleatorio.

Los cuatro ejemplos comparten demanda anual \(1200\), desviación \(100\),
plazo \(0.02\) años, preparación \(30\), mantenimiento \(2\), costo de
atraso \(20\) y adquisición \(5\). Todo faltante queda pendiente; el costo
de revisión es \(5\). El archivo de revisión continua hasta un nivel utiliza
un tamaño medio de pedido suministrado de \(200\).
\begin{center}
\begin{tabular}{lll}
\toprule
Archivo bajo \texttt{manual/examples/} & Política & Resultado principal \\
\midrule
\path{prob_qr.json} & Cantidad fija \((Q,r)\) & \(Q=194.866241,r=54.242691\) \\
\path{prob_sS.json} & Continua hasta nivel & \(r=54.095104,S=254.095104\) \\
\path{prob_periodica.json} & Intervalo fijo & \(T=0.170783,S=321.459989\) \\
\path{prob_opcional.json} & Periódica opcional & \(T=0.064550,s=163.951339\) \\
\bottomrule
\end{tabular}
\end{center}
La revisión opcional devuelve además \(S=353.687998\). En continua fija,
el stock de seguridad es \(30.242691\), el servicio \(0.983761\) y el
costo relevante \(450.217864\). En periódica fija, la protección incluye
revisión más plazo, \(0.190783\) años: confundir ambos intervalos reduce
artificialmente el stock necesario. Las fracciones de atraso y venta perdida,
sus costos y las distribuciones se validan en el núcleo.

\section{Procedimiento en macOS y trazabilidad}
Abra el JSON normalizado o seleccione una variante de inventarios en el creador
de modelos. Complete demanda, costos y los parámetros específicos de la
variante; valide antes de ejecutar. Mantenga visibles las unidades y revise las
advertencias. En los resultados, EOQ muestra costos y comparación con el lote
suministrado; descuentos muestra candidatos; vendedor de periódicos muestra
servicio y resultados de demanda; lotes muestra cantidades y costos por
período; revisión estocástica muestra niveles, stock de seguridad y costos.
El selector de análisis/JSON conserva el documento del núcleo para exportación
sin transcribir números de los gráficos. Las cantidades y fórmulas proceden de
QSBCore, tanto en CLI como en macOS.

Todos los archivos de este capítulo se validan con
\texttt{validate-inventory-json} y se resuelven con
\texttt{solve-inventory-json}; añadir \texttt{--backend validate} valida sin
resolver. El verificador \path{manual/verify_examples.py} comprueba los
resultados esperados de \path{manual/examples/manifest.json}; no necesita
material de referencia propietario.

\section{Ejercicios}
\begin{enumerate}
\item Compare pedidos enteros de 189 y 190 en EOQ y explique por qué conviene
comparar costos en vez de redondear a ciegas.
\item Duplique el plazo del ejemplo estocástico: compare cambios en demanda
protegida y stock de seguridad. \textquestiondown{}Se duplican ambos?
\item Baje el costo de atraso del ejemplo de lotes y compruebe si aparecen
inventarios negativos. Interprete el atraso operacionalmente.
\end{enumerate}
